\begin{ouverture}[calcul]{5}{Les fractions comme nombres}{Partie 1 · Nombres et calculs}
{Quotient · Égalités à trous · Demi-droite · Nombres mixtes}
\objectif{Comprendre une fraction comme le résultat exact d’une division.}
\objectif{Placer des fractions sur une demi-droite et graduer un segment.}
\objectif{Reconnaître une fraction entière, décimale ou non décimale.}
\prerequis{Fractions du cours moyen : parts égales, fractions plus grandes que 1.}
\end{ouverture}
\pagestyle{otMain}\essentiel\fontsize{11}{15}\selectfont
\phantomsection\label{cours-5-1}
\essSection{1}{De la part d’un tout au quotient}
\begin{definition}Si $b$ n’est pas nul, le quotient de $a$ par $b$ est le nombre qui, multiplié par $b$, donne $a$. On le note $\frac ab$ : c’est le résultat exact de la division de $a$ par $b$. Ainsi $4\times\frac34=3$.\end{definition}
\label{fig-cours-5-1-cours-quart-de-3}\figureGeo{cours-quart-de-3}

\begin{propriete}$\frac34$ désigne à la fois trois quarts d’une unité et le quart de 3 unités : les deux lectures donnent la même longueur. Une fraction est un nombre.\end{propriete}
\essSection{2}{Compléter une égalité multiplicative}
\begin{savoirfaire}[Égalités à trous]$7\times\square=3$ : le nombre cherché est $\frac37$, par définition du quotient. $\square\times4=3$ : c’est $\frac34$. $6\times\square=18$ : c’est $\frac{18}6$, soit 3.\end{savoirfaire}
\clearpage
\phantomsection\label{cours-5-2}
\essSection{3}{Repérer des fractions sur une demi-droite}
\begin{savoirfaire}[Placer une fraction]Pour placer $\frac74$, on partage l’unité en 4 et on compte 7 pas d’un quart : on dépasse 1 de trois quarts, car $\frac74=1+\frac34$.\end{savoirfaire}
\label{fig-cours-5-2-cours-droite-quarts}\figureGeo{cours-droite-quarts}

\begin{savoirfaire}[Graduer un segment]Pour partager un segment en 7 parts égales sans mesurer, on le pose sur un réseau de droites parallèles également espacées, appelé guide-âne : les droites le coupent en 7 parts égales.\end{savoirfaire}
\label{fig-cours-5-2-cours-guide-ane}\figureGeo{cours-guide-ane}

\essSection{4}{Fractions, décimaux et nombres mixtes}
\begin{propriete}Une fraction peut désigner un entier, $\frac{12}4=3$ ; un décimal non entier, $\frac34=0{,}75$ ; ou un nombre non décimal, comme $\frac13$ : $3\times0{,}3=0{,}9$ ; $3\times0{,}33=0{,}99$… on n’atteint jamais 1. À connaître : $\frac14=0{,}25$ ; $\frac12=0{,}5$ ; $\frac34=0{,}75$ ; $\frac32=1{,}5$ ; $\frac42=2$ ; $\frac52=2{,}5$.\end{propriete}

\clearpage\rubrique{Découverte}
\begin{blocLivre}{ex-5-1}\begin{exercice}[1]{Quotient exact}
Écrire sous forme de fraction le résultat exact de $5:8$, de $13:4$ et de $2:7$. Lesquels sont des nombres décimaux ?
\end{exercice}\end{blocLivre}
\begin{blocLivre}{ex-5-2}\begin{exercice}[1]{Égalités à trous}
Compléter : $9\times\square=4$ ; $\square\times5=3$ ; $11\times\square=11$ ; $6\times\square=18$.
\end{exercice}\end{blocLivre}
\begin{blocLivre}{ex-5-3}\begin{exercice}[1]{À connaître}
Donner l’écriture décimale de $\frac14$ ; $\frac34$ ; $\frac32$ ; $\frac52$ ; $\frac42$, puis une fraction égale à $0{,}5$.
\end{exercice}\end{blocLivre}
\begin{blocLivre}{ex-5-4}\begin{exercice}[1]{Lire sur la droite}
L’unité est partagée en quarts. Lire les abscisses des points A, B et C sous forme de fraction, puis de nombre mixte quand c’est utile.
\label{fig-ex-5-4}\figureGeo{fq-4}
\end{exercice}\end{blocLivre}
\begin{blocLivre}{ex-5-5}\begin{exercice}[1]{Le quart de 3}
La bande mesure 3 unités. Elle est partagée en 4 parts égales. Quelle est la longueur d’une part ? Comparer avec trois quarts d’unité.
\label{fig-ex-5-5}\figureGeo{fq-5}
\end{exercice}\end{blocLivre}

\rubrique{Entraînement}
\begin{blocLivre}{ex-5-6}\begin{exercice}[2]{Placer des fractions}
L’unité est partagée en sixièmes. Placer les points D, E et F d’abscisses $\frac56$ ; $\frac{11}6$ et $\frac43$.
\label{fig-ex-5-6}\figureGeo{fq-6}
\end{exercice}\end{blocLivre}
\begin{blocLivre}{ex-5-7}\begin{exercice}[2]{Entier, décimal ou non décimal}
Classer en trois groupes, entiers, décimaux non entiers et non décimaux : $\frac{15}5$ ; $\frac74$ ; $\frac23$ ; $\frac9{10}$ ; $\frac56$.
\end{exercice}\end{blocLivre}
\begin{blocLivre}{ex-5-8}\begin{exercice}[2]{Nombres mixtes}
Écrire $\frac{17}5$ et $\frac{23}6$ sous la forme d’un entier plus une fraction plus petite que 1, puis écrire $3+\frac27$ avec une seule fraction.
\end{exercice}\end{blocLivre}
\begin{blocLivre}{ex-5-9}\begin{exercice}[2]{Graduer un segment}
Le segment [AB] mesure 5 cm. À l’aide du réseau de droites parallèles également espacées, le partager en 3 parts égales. Quelle est la longueur exacte d’une part ? Donner son arrondi au centième.
\label{fig-ex-5-9}\figureGeo{fq-9}
\end{exercice}\end{blocLivre}
\begin{blocLivre}{ex-5-10}\begin{exercice}[2]{Encadrer entre deux entiers}
Encadrer chaque fraction entre deux entiers consécutifs : $\frac{17}3$ ; $\frac{47}{10}$ ; $\frac{22}7$.
\end{exercice}\end{blocLivre}

\rubrique{Problèmes}
\begin{blocLivre}{ex-5-11}\begin{exercice}[3]{Le partage des pizzas}
À une table, 3 pizzas sont partagées également entre 4 enfants. À une autre, 5 pizzas entre 8 enfants. Quelle part de pizza reçoit chaque enfant ? À quelle table mange-t-on le plus ?
\end{exercice}\end{blocLivre}
\begin{blocLivre}{ex-5-12}\begin{exercice}[3]{La ficelle}
Une ficelle de 2 m est coupée en 5 morceaux égaux, une autre de 2 m en 3 morceaux égaux. Donner la longueur exacte d’un morceau dans chaque cas, puis une écriture décimale exacte ou arrondie au centième.
\end{exercice}\end{blocLivre}
\begin{blocLivre}{ex-5-13}\begin{exercice}[3]{Le nombre caché}
Quel nombre faut-il multiplier par 8 pour obtenir 6 ? Donner son écriture fractionnaire la plus simple et son écriture décimale.
\end{exercice}\end{blocLivre}
\begin{blocLivre}{ex-5-14}\begin{exercice}[3]{Les nombres rompus}
Au Moyen Âge, on appelait les fractions des « nombres rompus ». Un champ de 7 arpents est partagé également entre 3 héritiers. Quelle part reçoit chacun ? L’écrire comme nombre mixte, puis donner son arrondi au dixième.
\end{exercice}\end{blocLivre}
\begin{blocLivre}{ex-5-15}\begin{exercice}[3]{Vrai ou faux}
Sacha affirme : « $\frac13=0{,}33$. » A-t-il raison ? Justifier par une multiplication.
\end{exercice}\end{blocLivre}

\clearpage
